\section{Conclusion}
\label{sec:conclusion}

We have introduced \lucid{}, the first large-scale multimodal driving dataset
specifically designed to support causal and counterfactual reasoning for LLM-based
autonomous driving. \lucid{} provides 12,847 richly annotated real-world driving
scenarios collected across six global cities, encompassing 847,293 language annotations
including 312,456 structured causal chain QA pairs and 198,327 counterfactual QA pairs
--- annotation types that are absent from all prior datasets.

Through a comprehensive sensor suite (8 cameras, 5 LiDARs, 5 radars, GPS/IMU, HD maps),
a rigorous three-stage human annotation pipeline involving 234 expert annotators, and
a systematic quality control protocol achieving $\kappa > 0.80$ across all annotation
types, \lucid{} establishes a high-quality and reproducible evaluation resource.

Our four benchmark tasks --- Causal Scene Understanding (CSU), Counterfactual
Reasoning (CFR), Instruction-following Navigation (IFN), and Risk-Aware Planning
(RAP) --- span the core capabilities required for explainable AD. Baseline evaluation
of six state-of-the-art vision-language models reveals that even the most capable
model (GPT-4V) achieves only 58.3\% accuracy on counterfactual reasoning (vs 33.3\%
random) and 42.7\% instruction success rate, confirming that current models fall
substantially short of human-level causal reasoning in driving contexts.

\lucid{} is publicly released under a Creative Commons BY-NC 4.0 licence.
We release the full dataset, annotation tool, evaluation server, baseline code,
and model checkpoints to enable reproducible research and to lower the barrier
to entry for the research community. We believe that \lucid{} will catalyse
progress toward LLM-based autonomous driving systems that are not only capable
but genuinely explainable, consistent, and trustworthy.

\section*{Acknowledgements}

The authors thank the 234 annotation team members across Shanghai, Munich, and London
for their diligent work, and the collection driver teams at Alpha, Beta, and Gamma fleets.
This work was supported in part by the National Natural Science Foundation of China
(Grant No.\ 62376158), the Swiss National Science Foundation (Grant No.\ 200021-212088),
and the MIT CSAIL Alliance Program.

\section*{Data Availability}

\lucid{} is publicly available at \url{https://bullaraodomathoti.github.io/lucid-drive/}. The dataset
is released under Creative Commons Attribution-NonCommercial 4.0 (CC BY-NC 4.0).
Evaluation is conducted via an online server at the same URL.

\section*{Declaration of Competing Interests}

The authors declare that they have no known competing financial interests or personal
relationships that could have appeared to influence the work reported in this paper.

\section*{Author Contributions}

\textbf{Wei Zhang}: Conceptualisation, Methodology, Data Collection, Writing --- Original Draft.
\textbf{Xinyu Liu}: Software, Annotation Pipeline, Formal Analysis.
\textbf{Sarah A. Johnson}: Methodology, Baseline Experiments, Writing --- Review \& Editing.
\textbf{Marc Hofmann}: Sensor System Design, Data Curation.
\textbf{Kai Chen}: Conceptualisation, Supervision, Funding Acquisition, Writing --- Review \& Editing.
